\documentclass[11pt,a4paper]{article}
\usepackage[utf8x]{inputenc}
\usepackage[T1]{fontenc}
\usepackage{color}
\usepackage{multicol}
\usepackage{array}%création tableaux
\usepackage{graphicx}%pour les photosimages
\usepackage{subcaption}
\usepackage{float}
% \graphicspath{{img/}}
\def\dbasepath{../../resources-figures-jerboa}
\def\dlogopath{../../logos}
\graphicspath{{\dbasepath/article2/}%
		{\dbasepath/figures/}%
		{\dbasepath/concepts/}%
		{\dbasepath/rules/}%
		{\dbasepath/persistent-naming/}%
		{\dlogopath}}

\usepackage{subcaption}
%\usepackage{gentium}

\usepackage{fancyhdr}
%\pagestyle{fancy}
%\renewcommand{\headrulewidth}{0pt}
%\renewcommand{\footrulewidth}{0pt}
%\setlength\headheight{80.0pt}
%\addtolength{\textheight}{-80.0pt}
%\chead{\includegraphics[scale=0.4]{Logo-template.PNG}}

\usepackage{mathptmx} % Use Times Font
\usepackage{geometry}
         \geometry{
         a4paper,
         left=14mm,
         top=20mm,
         right =15mm,
         bottom =15mm
         }
\pagestyle{fancy}
\renewcommand{\headrulewidth}{0pt}
\chead{\includegraphics[scale=9.]{logos_2024.pdf}}

\setlength\parindent{0em}

\begin{document}

\begin{center}
	{\Large \textbf{Réévaluation de modèle à base de règles de transformation de graphe}}
\end{center}

{\underline{M. Gaide}$^{1}$, S. Jean$^{1}$, D. Marcheix$^{1}$, X. Skapin$^{2}$, A. Arnould$^{2}$, H. Belhaouari$^{2}$}
\begin{flushright}
	\underline{Contact} : maxime.gaide@ensma.fr
\end{flushright}

\begin{flushleft}
	\scriptsize{
		$^1$	LIAS IDD, ISAE-ENSMA\\
		$^2$	XLIM ASALI, Université de Poitiers
	}
\end{flushleft}

\rule{1\linewidth}{3pt}

\textbf{Mots-clés} : Modélisation géométrique à base topologique, Règles de
transformation de graphe, Nommage Persistant, Réévaluation, Cartes généralisées

~

Nous étendons l'étude du problème de nommage[kripac,Capoyleas,bidarra] aux systèmes de modélisations basés
sur les règles de transformation de graphe[bohl2015modeling,terraz20093gmap]. %
Pour cela, nous proposons une méthode de nommage persistant s'appuyant sur les
cartes généralisées[1,2] et la formalisation des opérations par des règles de
transformation de graphe. %
Ces formalismes permettent la caractérisation unique et homogène des entités
topologiques en toutes dimensions. %

Afin de réévaluer des objets, la plupart des méthodes existantes nécessitent de
suivre la totalité, ou presque, des entités topologiques qui les composent. %
Par ailleurs, ces mêmes méthodes considèrent le problème du nommage seulement du
point de vue de la modification des paramètres dans une spécification
paramétriques. %

Avec notre méthode, non seulement nous abordons le nommage persistant dans le
cadre habituel de l'édition de paramètres, mais nous prenons aussi en compte
l'édition de la spécification paramétrique elle-même (ajout, suppression et
déplacement d'opération). %
De plus, nous utilisons des graphes orientés acycliques pour représenter
l'histoire des entités passées en paramètres dans les opérations jusqu'à leurs
origines. %

\begin{figure}[H]
	\centering
	\begin{subfigure}{.2\columnwidth}
		\centering
		\includegraphics[width=.6\textwidth]{short-1-evaluation-create-face}
		\caption{1-create(pos)}
		\label{fig:fr-create-face}
	\end{subfigure}
	\begin{subfigure}{.2\columnwidth}
		\centering
		\includegraphics[width=.6\textwidth]{short-2-evaluation-extrude-face}
		\caption{2-extrude($PN_1$, vec)}
		\label{fig:fr-extrude-face}
	\end{subfigure}
	\begin{subfigure}{.2\columnwidth}
		\centering
		\includegraphics[width=.6\textwidth]{short-3-evaluation-triangulate-face}
		\caption{3-triangulate($PN_2$)}
		\label{fig:fr-triangulate-face}
	\end{subfigure}
	\begin{subfigure}{.2\columnwidth}
		\centering
		\includegraphics[width=.6\textwidth]{short-4-evaluation-color-face}
		\caption{4-color($PN_3$)}
		\label{fig:fr-color-face}
	\end{subfigure}
	\caption{Evaluation du processus de construction d'un objet}
  \label{fig:fig-eval}
\end{figure}
\vspace*{-1em}
\begin{figure}[H]
	\centering
	\begin{subfigure}{.19\columnwidth}
		\centering
		\includegraphics[width=.6\textwidth]{short-1-reevaluation-create-face-darts}
		\caption{1-create(pos)}
		\label{fig:fr-reeval-create-face}
	\end{subfigure}
	\begin{subfigure}{.19\columnwidth}
		\centering
		\includegraphics[width=.6\textwidth]{short-1ADD-reevaluation-insert-vertex-darts}
		\caption{ADD1-insert(3)}
		\label{fig:fr-reeval-insert-vertex}
	\end{subfigure}
	\begin{subfigure}{.19\columnwidth}
		\centering
		\includegraphics[width=.6\textwidth]{short-2-reevaluation-extrude-face-darts}
		\caption{2-extrude($PN_1$, vec)}
		\label{fig:fr-reeval-extrude-face}
	\end{subfigure}
	\begin{subfigure}{.19\columnwidth}
		\centering
		\includegraphics[width=.6\textwidth]{short-3-reevaluation-triangulate-face-darts}
		\caption{3-triangulate($PN_2$)}
		\label{fig:fr-reeval-triangulate-face}
	\end{subfigure}
	\begin{subfigure}{.19\columnwidth}
		\centering
		\includegraphics[width=.6\textwidth]{short-4-reevaluation-color-face}
		\caption{4-color($PN_3$)}
		\label{fig:fr-reeval-color-face}
	\end{subfigure}
	\caption{Ajout d'une opération et réévaluation du processus de construction de l'objet en Figure~\ref{fig:fig-eval}}
  \label{fig:fig-reeval}
\end{figure}

% \begin{minipage}[b]{.45\textwidth}
% 	\begin{figure}[H]
% 		\centering
% 		\includegraphics[height=100pt]{pn7-short-nolevels}
% 		\caption{}
% 	\end{figure}
% \end{minipage}
% \begin{minipage}[b]{.45\textwidth}
% 	\begin{figure}[H]
% 		\centering
% 		\includegraphics[height=100pt]{pn7-reval-short-nolevels}
% 		\caption{}
% 	\end{figure}
% \end{minipage}

\vfill

\begin{flushright}
	\rule{0.85\linewidth}{3pt}
\end{flushright}

\textbf{Références (optionnelles)} {\color{red}(TNR, 11, Gras, justifié)}

{%
	\footnotesize %% Affecte tout ce qui suit, doit être contenu entre `{` et `}`
	$^{[1]}$ P.Nom 1, P.P. Nom 2, P.-P. Nom 3.Nom de la revue scientifique N°Vol (N°Sous Vol) (Année) Page initiale–Page finale.

	$^{[2]}$ P.Nom 1, P.P. Nom 2, P.-P. Nom 3.Nom de la revue scientifique N°Vol (N°Sous Vol) (Année) Page initiale–Page finale.

	$^{[3]}$ P.Nom 1, P.P. Nom 2, P.-P. Nom 3.Nom de la revue scientifique N°Vol (N°Sous Vol) (Année) Page initiale–Page finale.

}
\end{document}
